\documentclass[UTF8,11pt]{ctexart}
\usepackage[a4paper,margin=24mm]{geometry}
\usepackage{amsmath,amssymb,amsthm,mathtools,mathrsfs}
\usepackage[all]{xy}
\usepackage{xurl,hyperref,enumitem}
\usepackage{tikz}
\usetikzlibrary{arrows.meta,cd,decorations.markings,calc,patterns}
\hypersetup{hidelinks,breaklinks=true}
\setlength{\emergencystretch}{3em}
\allowdisplaybreaks
\newtheorem*{theorem}{Theorem}
\newtheorem*{thm}{Theorem}
\newtheorem*{lemma}{Lemma}
\newtheorem*{proposition}{Proposition}
\newtheorem*{prop}{Proposition}
\newtheorem*{corollary}{Corollary}
\newtheorem*{cor}{Corollary}
\newtheorem*{claim}{Claim}
\theoremstyle{definition}
\newtheorem*{definition}{Definition}
\newtheorem*{defn}{Definition}
\newtheorem*{remark}{Remark}
\newtheorem*{example}{Example}
\newcommand{\R}{\mathbb R}
\newcommand{\C}{\mathbb C}
\newcommand{\Z}{\mathbb Z}
\newcommand{\Q}{\mathbb Q}
\newcommand{\N}{\mathbb N}
\newcommand{\F}{\mathbb F}
\newcommand{\D}{\mathbb D}
\newcommand{\bB}{\mathbb B}
\newcommand{\bC}{\mathbb C}
\newcommand{\bR}{\mathbb R}
\newcommand{\bZ}{\mathbb Z}
\newcommand{\bN}{\mathbb N}
\newcommand{\bQ}{\mathbb Q}
\newcommand{\bD}{\mathbb D}
\newcommand{\bF}{\mathbb F}
\newcommand{\CP}{\mathbb{CP}}
\newcommand{\RP}{\mathbb{RP}}
\newcommand{\sO}{\mathscr O}
\newcommand{\sI}{\mathscr I}
\newcommand{\sF}{\mathscr F}
\newcommand{\sG}{\mathscr G}
\newcommand{\sH}{\mathscr H}
\newcommand{\sR}{\mathscr R}
\newcommand{\sM}{\mathscr M}
\newcommand{\sV}{\mathscr V}
\newcommand{\sU}{\mathscr U}
\DeclareMathOperator{\Hom}{Hom}
\DeclareMathOperator{\Ext}{Ext}
\DeclareMathOperator{\Tor}{Tor}
\DeclareMathOperator{\Spec}{Spec}
\DeclareMathOperator{\Proj}{Proj}
\DeclareMathOperator{\supp}{supp}
\DeclareMathOperator{\im}{im}
\DeclareMathOperator{\id}{id}
\DeclareMathOperator{\Tr}{Tr}
\DeclareMathOperator{\tr}{tr}
\DeclareMathOperator{\rank}{rank}
\DeclareMathOperator{\ord}{ord}
\DeclareMathOperator{\dist}{dist}
\DeclareMathOperator{\End}{End}
\DeclareMathOperator{\Aut}{Aut}
\DeclareMathOperator{\Gal}{Gal}
\DeclareMathOperator{\vol}{vol}
\DeclareMathOperator{\Cl}{Cl}
\DeclareMathOperator{\coker}{coker}
\DeclareMathOperator{\codim}{codim}
\DeclareMathOperator{\divergence}{div}
\DeclareMathOperator{\Ric}{Ric}
\DeclareMathOperator{\grad}{grad}
\DeclareMathOperator{\Hess}{Hess}
\newcommand{\abs}[1]{\left\lvert#1\right\rvert}
\newcommand{\sabs}[1]{\lvert#1\rvert}
\newcommand{\babs}[1]{\bigl\lvert#1\bigr\rvert}
\newcommand{\Babs}[1]{\Bigl\lvert#1\Bigr\rvert}
\newcommand{\bbabs}[1]{\biggl\lvert#1\biggr\rvert}
\newcommand{\BBabs}[1]{\Biggl\lvert#1\Biggr\rvert}
\newcommand{\norm}[1]{\left\lVert#1\right\rVert}
\newcommand{\snorm}[1]{\lVert#1\rVert}
\newcommand{\bnorm}[1]{\bigl\lVert#1\bigr\rVert}
\newcommand{\Bnorm}[1]{\Bigl\lVert#1\Bigr\rVert}
\newcommand{\bbnorm}[1]{\biggl\lVert#1\biggr\rVert}
\newcommand{\BBnorm}[1]{\Biggl\lVert#1\Biggr\rVert}
\newcommand{\widebar}[1]{\overline{#1}}
\newcommand{\nicefrac}[2]{\frac{#1}{#2}}
\newcommand{\linnprod}[2]{\langle#1,#2\rangle}
\newcommand{\myindex}[1]{#1}
\newcommand{\glsadd}[1]{}
\newcommand{\avoidbreak}{}
\newcommand{\sourceref}[1]{\textnormal{[原文：\nolinkurl{#1}]}}
\newcommand{\thmref}[1]{\sourceref{#1}}
\newcommand{\lemmaref}[1]{\sourceref{#1}}
\newcommand{\propref}[1]{\sourceref{#1}}
\newcommand{\corref}[1]{\sourceref{#1}}
\newcommand{\defnref}[1]{\sourceref{#1}}
\newcommand{\exampleref}[1]{\sourceref{#1}}
\newcommand{\exerciseref}[1]{\sourceref{#1}}
\newcommand{\figureref}[1]{\sourceref{#1}}
\newcommand{\sectionref}[1]{\sourceref{#1}}
\newcommand{\appendixref}[1]{\sourceref{#1}}
\newcommand{\stacksref}[2]{\href{https://stacks.math.columbia.edu/tag/#1}{#2 [#1]}}

\begin{document}
\section*{039\quad 拟有限分离态射的Zariski主定理}
\subsection*{来源与计数目标}
The Stacks Project Authors, \emph{The Stacks Project}。Section 37.43，Lemmas 37.43.1--37.43.3；另收 Lemma 37.41.4。使用实际取得的网页数学 TeX 正文，获取日期：2026-09-28；未取得网页构建的固定提交号，不编造版本。
\par\url{https://stacks.math.columbia.edu/tag/05K0}
\par\url{https://stacks.math.columbia.edu/tag/03GW}
\par\url{https://stacks.math.columbia.edu/tag/02LQ}
\par\url{https://stacks.math.columbia.edu/tag/02LN}

\par\textbf{本文件只计一个问题：}拟有限分离态射在拟紧拟分离底概形上分解为拟紧开浸入与有限态射。
\subsection*{作者命题与人类证明}

\begin{lemma}[37.41.4: finite pieces after an elementary étale neighbourhood]
Let $f:X\to S$ be a morphism of schemes. Let $s\in S$. Let $x_1,\ldots,x_n\in X_s$. Assume that $f$ is locally of finite type and separated, and that $x_1,\ldots,x_n$ are pairwise distinct isolated points of $X_s$. Then there exists an elementary étale neighbourhood $(U,u)\to(S,s)$ and a decomposition
\[
U\times_S X=W\amalg V_1\amalg\cdots\amalg V_n
\]
into open and closed subschemes such that the morphisms $V_i\to U$ are finite, the fibres of $V_i\to U$ over $u$ are singletons $\{v_i\}$, each $v_i$ maps to $x_i$ with $\kappa(x_i)=\kappa(v_i)$, and the fibre of $W\to U$ over $u$ contains no points mapping to any of the $x_i$.
\end{lemma}
\begin{proof}
Choose $(U,u)\to(S,s)$ and $V_i\subset X_U$ as in Lemma 37.41.2. Since $X_U\to U$ is separated (Schemes, Lemma 26.21.12) and $V_i\to U$ is finite hence proper (Morphisms, Lemma 29.45.11) we see that $V_i\subset X_U$ is closed by Morphisms, Lemma 29.42.7. Hence $V_i\cap V_j$ is a closed subset of $V_i$ which does not contain $v_i$. Hence the image of $V_i\cap V_j$ in $U$ is a closed set (because $V_i\to U$ proper) not containing $u$. After shrinking $U$ we may therefore assume that $V_i\cap V_j=\emptyset$ for all $i,j$. This gives the decomposition as in the lemma.
\end{proof}

\begin{lemma}[37.43.1]
Let $f:X\to S$ be a morphism of schemes. Assume $f$ is of finite type and separated. Let $S'$ be the normalization of $S$ in $X$, see Morphisms, Definition 29.54.3. Picture:
\[
\xymatrix{X\ar[rd]_f\ar[rr]^{f'}&&S'\ar[ld]^\nu\\&S&}
\]
Then there exists an open subscheme $U'\subset S'$ such that
\begin{enumerate}
\item $(f')^{-1}(U')\to U'$ is an isomorphism, and
\item $(f')^{-1}(U')\subset X$ is the set of points at which $f$ is quasi-finite.
\end{enumerate}
\end{lemma}
\begin{proof}
By Morphisms, Lemma 29.57.2 the subset $U\subset X$ of points where $f$ is quasi-finite is open. The lemma is equivalent to
\begin{enumerate}[label=(\alph*)]
\item $U'=f'(U)\subset S'$ is open,
\item $U=(f')^{-1}(U')$, and
\item $U\to U'$ is an isomorphism.
\end{enumerate}
Let $x\in U$ be arbitrary. We claim there exists an open neighbourhood $f'(x)\in V\subset S'$ such that $(f')^{-1}V\to V$ is an isomorphism. We first prove the claim implies the lemma. Namely, then $(f')^{-1}V\cong V$ is both locally of finite type over $S$ (as an open subscheme of $X$) and for $v\in V$ the residue field extension $\kappa(v)/\kappa(\nu(v))$ is algebraic (as $V\subset S'$ and $S'$ is integral over $S$). Hence the fibres of $V\to S$ are discrete (Morphisms, Lemma 29.21.2) and $(f')^{-1}V\to S$ is locally quasi-finite (Morphisms, Lemma 29.21.8). This implies $(f')^{-1}V\subset U$ and $V\subset U'$. Since $x$ was arbitrary we see that (a), (b), and (c) are true.

Let $s=f(x)$. Let $(T,t)\to(S,s)$ be an elementary étale neighbourhood. Denote by a subscript ${}_T$ the base change to $T$. Let $y=(x,t)\in X_T$ be the unique point in the fibre $X_t$ lying over $x$. Note that $U_T\subset X_T$ is the set of points where $f_T$ is quasi-finite, see Morphisms, Lemma 29.21.13. Note that
\[
X_T\xrightarrow{f'_T}S'_T\xrightarrow{\nu_T}T
\]
is the normalization of $T$ in $X_T$, see Lemma 37.19.2. Suppose that the claim holds for $y\in U_T\subset X_T\to S'_T\to T$, i.e., suppose that we can find an open neighbourhood $f'_T(y)\in V'\subset S'_T$ such that $(f'_T)^{-1}V'\to V'$ is an isomorphism. The morphism $S'_T\to S'$ is étale hence the image $V\subset S'$ of $V'$ is open. Observe that $f'(x)\in V$ as $f'_T(y)\in V'$. Observe that
\[
\xymatrix{(f'_T)^{-1}V'\ar[r]\ar[d]&(f')^{-1}(V)\ar[d]\\V'\ar[r]&V}
\]
is a fibre square (as $S'_T\times_{S'}X=X_T$). Since the left vertical arrow is an isomorphism and $\{V'\to V\}$ is a étale covering, we conclude that the right vertical arrow is an isomorphism by Descent, Lemma 35.23.19. In other words, the claim holds for $x\in U\subset X\to S'\to S$.

By the result of the previous paragraph we may replace $S$ by an elementary étale neighbourhood of $s=f(x)$ in order to prove the claim. Thus we may assume there is a decomposition
\[X=V\amalg W\]
into open and closed subschemes where $V\to S$ is finite and $x\in V$, see Lemma 37.41.4. Since $X$ is a disjoint union of $V$ and $W$ over $S$ and since $V\to S$ is finite we see that the normalization of $S$ in $X$ is the morphism
\[X=V\amalg W\longrightarrow V\amalg W'\longrightarrow S\]
where $W'$ is the normalization of $S$ in $W$, see Morphisms, Lemmas 29.54.10, 29.45.4, and 29.54.12. The claim follows and we win.
\end{proof}

\begin{lemma}[37.43.2]
Let $f:X\to S$ be a morphism of schemes. Assume $f$ is quasi-finite and separated. Let $S'$ be the normalization of $S$ in $X$. Then $f':X\to S'$ is a quasi-compact open immersion and $\nu:S'\to S$ is integral. In particular $f$ is quasi-affine.
\end{lemma}
\begin{proof}
This follows from Lemma 37.43.1. Namely, by that lemma there exists an open subscheme $U'\subset S'$ such that $(f')^{-1}(U')=X$ and $X\to U'$ is an isomorphism. In other words, $f'$ is an open immersion. Note that $f'$ is quasi-compact as $f$ is quasi-compact and $\nu:S'\to S$ is separated (Schemes, Lemma 26.21.14). It follows that $f$ is quasi-affine by Morphisms, Lemma 29.13.3.
\end{proof}

\begin{lemma}[37.43.3: Zariski's Main Theorem]
Let $f:X\to S$ be a morphism of schemes. Assume $f$ is quasi-finite and separated and assume that $S$ is quasi-compact and quasi-separated. Then there exists a factorization
\[
\xymatrix{X\ar[rd]_f\ar[rr]_j&&T\ar[ld]^\pi\\&S&}
\]
where $j$ is a quasi-compact open immersion and $\pi$ is finite.
\end{lemma}
\begin{proof}
Let $X\to S'\to S$ be as in the conclusion of Lemma 37.43.2. By Properties, Lemma 28.23.13 we can write
\[
\nu_*\mathcal O_{S'}=\mathop{\mathrm{colim}}_{i\in I}\mathcal A_i
\]
as a directed colimit of finite quasi-coherent $\mathcal O_S$-algebras $\mathcal A_i\subset\nu_*\mathcal O_{S'}$. Then $\pi_i:T_i=\underline{\Spec}_S(\mathcal A_i)\to S$ is a finite morphism for each $i$. Note that the transition morphisms $T_{i'}\to T_i$ are affine and that $S'=\mathop{\mathrm{lim}}T_i$. By Limits, Lemma 32.4.11 there exists an $i$ and a quasi-compact open $U_i\subset T_i$ whose inverse image in $S'$ equals $f'(X)$. For $i'\geq i$ let $U_{i'}$ be the inverse image of $U_i$ in $T_{i'}$. Then $X\cong f'(X)=\mathop{\mathrm{lim}}_{i'\geq i}U_{i'}$, see Limits, Lemma 32.2.2. By Limits, Lemma 32.4.16 we see that $X\to U_{i'}$ is a closed immersion for some $i'\geq i$. (In fact $X\cong U_{i'}$ for sufficiently large $i'$ but we don't need this.) Hence $X\to T_{i'}$ is an immersion. By Morphisms, Lemma 29.3.2 we can factor this as $X\to T\to T_{i'}$ where the first arrow is an open immersion and the second a closed immersion. Thus we win.
\end{proof}

\subsection*{整理说明（非作者正文）}
允许的先修：相对正规化与 étale 基变换相容；拟有限态射的初等 étale 局部有限邻域定理 37.41.2；同构的 étale 下降；仿射逆极限的有限阶段逼近（32.4.11、32.4.16）及正文所列概形基本性质。本题在这些工具之上仍需把局部有限分支识别为相对正规化中的开集，下降该识别，再从积分解构造有限分解，相关核心论证均在正文。不是把 Zariski 主定理本身当作先修；不将支持引理及代数版再计数。保留交换图和作者论证顺序，外部引用采用原章节号；保存所用网页 TeX 摘录，不是整章原件。
\subsection*{可修改的简短标注（非作者正文）}
$c$：拟有限态射的有限分解。

$a$：相对正规化；初等 étale 邻域；有限分支；同构下降；拟凝聚代数的有向余极限；仿射逆极限有限阶段逼近。

\texttt{cross\_theory\_status = yes}。在 37.43.1 中以 étale 局部有限分支控制正规化的开子概形，再通过下降返回原底；37.43.3 将积分解的代数余极限转成有限态射，实现所需的全局分解。

全部标注均为非穷尽、可修改初稿；未标注不表示未使用。
\subsection*{许可与修改记录}
原数学正文作者：The Stacks Project Authors。按 GNU Free Documentation License 1.2 或后续版本复用；无不变章节、封面文字或封底文字。许可全文位于 \texttt{sources/stacks/GFDL-1.2.txt}。本条整理部分沿用同一许可。2026-09-28：助手摘录、独立排版并加入来源、范围说明与初稿标注；原作者不对本次整理背书。
\end{document}
